\documentclass[12pt,a4paper]{article}

\usepackage[T2A]{fontenc}
\usepackage[utf8]{inputenc}
\usepackage[russian]{babel}
\usepackage{amsmath,amssymb,amsthm,mathtools}
\usepackage{array}
\usepackage{enumitem}
\usepackage{listings}
\usepackage{xcolor}
\usepackage{geometry}

\geometry{left=25mm,right=20mm,top=20mm,bottom=20mm}

\lstset{
  basicstyle=\ttfamily\small,
  keywordstyle=\bfseries\color{blue!60!black},
  commentstyle=\itshape\color{green!40!black},
  stringstyle=\color{red!60!black},
  numbers=left,
  numberstyle=\tiny,
  stepnumber=1,
  numbersep=8pt,
  frame=single,
  breaklines=true,
  tabsize=2,
  showstringspaces=false,
  literate=
    {А}{{\CYRA}}1
    {Б}{{\CYRB}}1
    {В}{{\CYRV}}1
    {Г}{{\CYRG}}1
    {Д}{{\CYRD}}1
    {Е}{{\CYRE}}1
    {Ё}{{\CYRYO}}1
    {Ж}{{\CYRZH}}1
    {З}{{\CYRZ}}1
    {И}{{\CYRI}}1
    {Й}{{\CYRISHRT}}1
    {К}{{\CYRK}}1
    {Л}{{\CYRL}}1
    {М}{{\CYRM}}1
    {Н}{{\CYRN}}1
    {О}{{\CYRO}}1
    {П}{{\CYRP}}1
    {Р}{{\CYRR}}1
    {С}{{\CYRS}}1
    {Т}{{\CYRT}}1
    {У}{{\CYRU}}1
    {Ф}{{\CYRF}}1
    {Х}{{\CYRH}}1
    {Ц}{{\CYRC}}1
    {Ч}{{\CYRCH}}1
    {Ш}{{\CYRSH}}1
    {Щ}{{\CYRSHCH}}1
    {Ъ}{{\CYRHRDSN}}1
    {Ы}{{\CYRERY}}1
    {Ь}{{\CYRSFTSN}}1
    {Э}{{\CYREREV}}1
    {Ю}{{\CYRYU}}1
    {Я}{{\CYRYA}}1
    {а}{{\cyra}}1
    {б}{{\cyrb}}1
    {в}{{\cyrv}}1
    {г}{{\cyrg}}1
    {д}{{\cyrd}}1
    {е}{{\cyre}}1
    {ё}{{\cyryo}}1
    {ж}{{\cyrzh}}1
    {з}{{\cyrz}}1
    {и}{{\cyri}}1
    {й}{{\cyrishrt}}1
    {к}{{\cyrk}}1
    {л}{{\cyrl}}1
    {м}{{\cyrm}}1
    {н}{{\cyrn}}1
    {о}{{\cyro}}1
    {п}{{\cyrp}}1
    {р}{{\cyrr}}1
    {с}{{\cyrs}}1
    {т}{{\cyrt}}1
    {у}{{\cyru}}1
    {ф}{{\cyrf}}1
    {х}{{\cyrh}}1
    {ц}{{\cyrc}}1
    {ч}{{\cyrch}}1
    {ш}{{\cyrsh}}1
    {щ}{{\cyrshch}}1
    {ъ}{{\cyrhrdsn}}1
    {ы}{{\cyrery}}1
    {ь}{{\cyrsftsn}}1
    {э}{{\cyrerev}}1
    {ю}{{\cyryu}}1
    {я}{{\cyrya}}1
}

\theoremstyle{definition}
\newtheorem{definition}{Определение}
\newtheorem{example}{Пример}
\newtheorem{remark}{Замечание}
\newtheorem{statement}{Утверждение}

\theoremstyle{plain}
\newtheorem{theorem}{Теорема}
\newtheorem{proposition}{Утверждение}

\title{Билет №54 \\ \small Автоматическое построение лексических и синтаксических анализаторов по формальным описаниям грамматик. Инструменты генерации лексических и синтаксических анализаторов: lex, yacc, ANTLR, парсер-комбинаторы. (СпБГУ)}
\author{Сериков Владислав Александрович}
\date{16 мая 2026}

\begin{document}

\maketitle
\tableofcontents
\newpage

\section{Автоматическое построение лексических и синтаксических анализаторов}

\subsection{Общая схема фронтенда компилятора}

\textbf{Лексический анализатор} разбивает входной поток символов на поток токенов:
\[
\text{символы} \longrightarrow \text{токены}.
\]
Например, строка
\[
\texttt{x = 12 + y}
\]
может быть преобразована в поток:
\[
\langle ID(x), ASSIGN, NUM(12), PLUS, ID(y), EOF\rangle.
\]

\textbf{Синтаксический анализатор} проверяет, соответствует ли поток токенов грамматике языка, и строит дерево разбора или абстрактное синтаксическое дерево:
\[
\text{токены} \longrightarrow \text{дерево разбора / AST}.
\]

Типичная архитектура:
\[
\boxed{
\text{исходный текст}
\to
\text{лексер}
\to
\text{поток токенов}
\to
\text{парсер}
\to
\text{AST}
}
\]

Автоматическое построение анализаторов основано на том, что лексическая структура языков обычно описывается \textbf{регулярными выражениями}, а синтаксическая структура --- \textbf{контекстно-свободными грамматиками}.

\subsection{Алфавиты, слова, языки}

\begin{definition}
\textbf{Алфавит} $\Sigma$ --- конечное непустое множество символов.
\end{definition}

\begin{definition}
\textbf{Слово} над алфавитом $\Sigma$ --- конечная последовательность символов из $\Sigma$.
\end{definition}

Пустое слово обозначается $\varepsilon$.

Множество всех слов над $\Sigma$ обозначается $\Sigma^*$.

\begin{definition}
\textbf{Язык} над алфавитом $\Sigma$ --- произвольное подмножество $\Sigma^*$.
\end{definition}

Например, если
\[
\Sigma = \{a,b\},
\]
то
\[
\Sigma^* = \{\varepsilon, a, b, aa, ab, ba, bb, aaa, \ldots\}.
\]

\subsection{Регулярные языки и регулярные выражения}

Лексические классы обычно описываются регулярными выражениями.

\begin{definition}
Регулярные выражения над алфавитом $\Sigma$ определяются индуктивно:
\begin{enumerate}
    \item $\varnothing$ --- регулярное выражение;
    \item $\varepsilon$ --- регулярное выражение;
    \item если $a \in \Sigma$, то $a$ --- регулярное выражение;
    \item если $r$ и $s$ --- регулярные выражения, то $(r|s)$, $(rs)$ и $(r^*)$ также регулярные выражения.
\end{enumerate}
\end{definition}

Семантика:
\[
L(\varnothing)=\varnothing,
\]
\[
L(\varepsilon)=\{\varepsilon\},
\]
\[
L(a)=\{a\},
\]
\[
L(r|s)=L(r)\cup L(s),
\]
\[
L(rs)=\{xy \mid x\in L(r),\ y\in L(s)\},
\]
\[
L(r^*)=\bigcup_{k\geq 0} L(r)^k.
\]

Примеры:
\[
[0-9]^+
\]
задаёт непустые последовательности цифр.

\[
[a-zA-Z\_][a-zA-Z0-9\_]*
\]
задаёт идентификаторы в стиле языков C.

\subsection{Конечные автоматы}

\begin{definition}
\textbf{Детерминированный конечный автомат} --- пятёрка
\[
D = (Q,\Sigma,\delta,q_0,F),
\]
где:
\begin{itemize}
    \item $Q$ --- конечное множество состояний;
    \item $\Sigma$ --- входной алфавит;
    \item $\delta: Q \times \Sigma \to Q$ --- функция переходов;
    \item $q_0\in Q$ --- начальное состояние;
    \item $F\subseteq Q$ --- множество допускающих состояний.
\end{itemize}
\end{definition}

\begin{definition}
\textbf{Недетерминированный конечный автомат} --- пятёрка
\[
N=(Q,\Sigma,\Delta,q_0,F),
\]
где
\[
\Delta: Q\times(\Sigma\cup\{\varepsilon\})\to 2^Q.
\]
\end{definition}

Недетерминированный автомат может:
\begin{itemize}
    \item иметь несколько переходов по одному символу;
    \item иметь $\varepsilon$-переходы, не потребляющие символ входа.
\end{itemize}

\subsection{Теорема Клини}

\begin{theorem}[Клини]
Класс языков, задаваемых регулярными выражениями, совпадает с классом языков, распознаваемых конечными автоматами.
\end{theorem}

\begin{proof}
Докажем две стороны.

\textbf{1. Всякий регулярный язык распознаётся конечным автоматом.}

Докажем индукцией по построению регулярного выражения.

\begin{enumerate}
    \item Для $\varnothing$ строится автомат без допускающих путей.
    \item Для $\varepsilon$ строится автомат:
    \[
    q_0 \xrightarrow{\varepsilon} q_f.
    \]
    \item Для символа $a$ строится автомат:
    \[
    q_0 \xrightarrow{a} q_f.
    \]
    \item Для объединения $r|s$ берём автоматы для $r$ и $s$, добавляем новое начальное состояние и $\varepsilon$-переходы к начальным состояниям обоих автоматов. Новое допускающее состояние достижимо из допускающих состояний обоих автоматов.
    \item Для конкатенации $rs$ соединяем каждое допускающее состояние автомата для $r$ с начальным состоянием автомата для $s$ $\varepsilon$-переходом.
    \item Для итерации $r^*$ добавляем возможность:
    \begin{itemize}
        \item принять пустое слово;
        \item после распознавания одного экземпляра $r$ вернуться к началу автомата для $r$.
    \end{itemize}
\end{enumerate}

Таким образом, по любому регулярному выражению можно построить НКА.

\textbf{2. Всякий язык, распознаваемый конечным автоматом, задаётся регулярным выражением.}

Рассмотрим конечный автомат с состояниями
\[
q_1,\ldots,q_n.
\]
Для каждой пары состояний $q_i,q_j$ и числа $k$ определим регулярное выражение
\[
R_{ij}^{(k)},
\]
задающее множество всех слов, переводящих автомат из $q_i$ в $q_j$ с использованием в качестве промежуточных только состояний из множества
\[
\{q_1,\ldots,q_k\}.
\]

База:
\[
R_{ij}^{(0)}
\]
задаёт все прямые переходы из $q_i$ в $q_j$, а также $\varepsilon$, если $i=j$.

Переход:
\[
R_{ij}^{(k)}
=
R_{ij}^{(k-1)}
\mid
R_{ik}^{(k-1)}
\left(R_{kk}^{(k-1)}\right)^*
R_{kj}^{(k-1)}.
\]

Идея формулы:
\begin{itemize}
    \item либо путь из $q_i$ в $q_j$ не проходит через $q_k$;
    \item либо впервые входит в $q_k$, затем некоторое число раз делает циклы через $q_k$, затем выходит из $q_k$ в $q_j$.
\end{itemize}

Итоговое регулярное выражение для языка автомата:
\[
\bigcup_{q_f\in F} R_{0f}^{(n)}.
\]

Следовательно, всякий автоматный язык является регулярным.

Обе стороны доказаны, значит классы совпадают.
\end{proof}

\subsection{Построение лексического анализатора}

Генератор лексических анализаторов принимает набор правил вида:
\[
\text{регулярное выражение} \quad \{ \text{действие} \}.
\]

Например:
\begin{verbatim}
[0-9]+                 { return NUM; }
[a-zA-Z_][a-zA-Z0-9_]* { return ID; }
[ \t\n]+               { skip; }
\end{verbatim}

Общая схема генерации:
\[
\boxed{
\text{регулярные выражения}
\to
\text{НКА}
\to
\text{ДКА}
\to
\text{минимизированный ДКА}
\to
\text{программа-лексер}
}
\]

\subsection{Алгоритм Томпсона: регулярное выражение $\to$ НКА}

Алгоритм Томпсона строит НКА с $\varepsilon$-переходами по структуре регулярного выражения.

\subsubsection*{Базовые случаи}

Для символа $a$:
\[
q_0 \xrightarrow{a} q_1.
\]

Для $\varepsilon$:
\[
q_0 \xrightarrow{\varepsilon} q_1.
\]

Для $\varnothing$ строится фрагмент без пути от начального состояния к конечному.

\subsubsection*{Объединение}

Для $r|s$:
\[
\begin{array}{ccccc}
& & q_r & \cdots & f_r \\
q_0 & \xrightarrow{\varepsilon} & & & \xrightarrow{\varepsilon} q_f\\
& & q_s & \cdots & f_s
\end{array}
\]

\subsubsection*{Конкатенация}

Для $rs$ конечное состояние автомата для $r$ соединяется $\varepsilon$-переходом с начальным состоянием автомата для $s$:
\[
q_r \cdots f_r \xrightarrow{\varepsilon} q_s \cdots f_s.
\]

\subsubsection*{Звезда Клини}

Для $r^*$:
\[
q_0 \xrightarrow{\varepsilon} q_f,
\]
\[
q_0 \xrightarrow{\varepsilon} q_r,
\]
\[
f_r \xrightarrow{\varepsilon} q_r,
\]
\[
f_r \xrightarrow{\varepsilon} q_f.
\]

\subsubsection*{Пример}

Построим НКА для выражения:
\[
(a|b)^*abb.
\]

Сначала строится НКА для $a|b$, затем добавляется звезда Клини, затем конкатенируются автоматы для $a$, $b$, $b$.

Этот автомат распознаёт все слова над $\{a,b\}$, заканчивающиеся на $abb$:
\[
abb,\ aabb,\ babb,\ aaabb,\ ababb,\ldots
\]

\subsection{Детерминизация НКА: алгоритм подмножеств}

\begin{definition}
$\varepsilon$-замыкание множества состояний $S$, обозначаемое
\[
\varepsilon\text{-closure}(S),
\]
--- множество всех состояний, достижимых из состояний $S$ по нулю или более $\varepsilon$-переходам.
\end{definition}

\subsubsection*{Алгоритм подмножеств}

Вход:
\[
N=(Q,\Sigma,\Delta,q_0,F).
\]

Выход:
\[
D=(Q_D,\Sigma,\delta_D,S_0,F_D).
\]

Шаги:
\begin{enumerate}
    \item Начальное состояние ДКА:
    \[
    S_0=\varepsilon\text{-closure}(\{q_0\}).
    \]
    \item Пока есть необработанное состояние $S\subseteq Q$:
    \begin{enumerate}
        \item для каждого символа $a\in\Sigma$ вычислить:
        \[
        T=\varepsilon\text{-closure}
        \left(
        \bigcup_{q\in S}\Delta(q,a)
        \right);
        \]
        \item добавить переход:
        \[
        \delta_D(S,a)=T;
        \]
        \item если $T$ ещё не встречалось, добавить его в множество состояний ДКА.
    \end{enumerate}
    \item Допускающие состояния ДКА:
    \[
    F_D=\{S\subseteq Q \mid S\cap F\neq \varnothing\}.
    \]
\end{enumerate}

\subsubsection*{Минимальный пример}

Пусть НКА распознаёт $a|b$:
\[
q_0 \xrightarrow{\varepsilon} q_1,\quad q_1\xrightarrow{a}q_2,
\]
\[
q_0 \xrightarrow{\varepsilon} q_3,\quad q_3\xrightarrow{b}q_4.
\]

Тогда:
\[
S_0=\{q_0,q_1,q_3\}.
\]

По символу $a$:
\[
\{q_0,q_1,q_3\}\xrightarrow{a}\{q_2\}.
\]

По символу $b$:
\[
\{q_0,q_1,q_3\}\xrightarrow{b}\{q_4\}.
\]

Состояния $\{q_2\}$ и $\{q_4\}$ являются допускающими.

\subsection{Правило максимального соответствия в лексерах}

Если несколько токенов могут быть распознаны из текущей позиции, лексер обычно применяет два правила:

\begin{enumerate}
    \item выбирается самый длинный возможный префикс;
    \item если несколько правил дают одинаково длинный префикс, выбирается правило, записанное раньше.
\end{enumerate}

Например, для входа:
\[
\texttt{ifx}
\]
и правил:
\[
\texttt{if} \to IF,
\]
\[
[a-z]+ \to ID,
\]
будет выбран токен:
\[
ID(\texttt{ifx}),
\]
потому что он длиннее, чем \texttt{if}.

Для входа:
\[
\texttt{if}
\]
оба правила подходят, но выбирается первое:
\[
IF.
\]

\subsection{Контекстно-свободные грамматики}

\begin{definition}
\textbf{Контекстно-свободная грамматика} --- четвёрка
\[
G=(N,\Sigma,P,S),
\]
где:
\begin{itemize}
    \item $N$ --- конечное множество нетерминалов;
    \item $\Sigma$ --- конечное множество терминалов;
    \item $P$ --- конечное множество продукций вида
    \[
    A\to \alpha,
    \]
    где $A\in N$, $\alpha\in (N\cup\Sigma)^*$;
    \item $S\in N$ --- стартовый символ.
\end{itemize}
\end{definition}

Пример грамматики арифметических выражений:
\[
E \to E + T \mid T,
\]
\[
T \to T * F \mid F,
\]
\[
F \to (E) \mid id.
\]

Эта грамматика задаёт выражения с приоритетом умножения над сложением.

\subsection{Дерево вывода и левосторонний вывод}

Для грамматики:
\[
E \to E+T \mid T,
\]
\[
T \to id
\]
слово
\[
id+id
\]
имеет левосторонний вывод:
\[
E \Rightarrow E+T \Rightarrow T+T \Rightarrow id+T \Rightarrow id+id.
\]

Дерево разбора:

\[
\begin{array}{c}
E\\
/\ \backslash\\
E\quad +\quad T\\
|\quad\quad\quad |\\
T\quad\quad\quad id\\
|\\
id
\end{array}
\]

\subsection{Неоднозначность грамматики}

\begin{definition}
Грамматика называется \textbf{неоднозначной}, если существует слово, имеющее два различных дерева разбора.
\end{definition}

Классический пример:
\[
E \to E+E \mid E*E \mid id.
\]

Слово:
\[
id+id*id
\]
имеет два разбора:
\[
(id+id)*id
\]
и
\[
id+(id*id).
\]

Для устранения неоднозначности вводят уровни приоритета:
\[
E \to E+T \mid T,
\]
\[
T \to T*F \mid F,
\]
\[
F \to id \mid (E).
\]

\subsection{FIRST и FOLLOW}

Множества FIRST и FOLLOW используются при построении LL-анализаторов.

\begin{definition}
\[
FIRST(\alpha)
\]
--- множество терминалов, с которых может начинаться строка, выводимая из $\alpha$.
Если $\alpha\Rightarrow^*\varepsilon$, то $\varepsilon\in FIRST(\alpha)$.
\end{definition}

\begin{definition}
\[
FOLLOW(A)
\]
--- множество терминалов, которые могут непосредственно следовать за нетерминалом $A$ в некоторой сентенциальной форме.
Для стартового символа $S$ всегда:
\[
EOF\in FOLLOW(S).
\]
\end{definition}

\subsubsection*{Пример}

Пусть:
\[
S\to AB,
\]
\[
A\to a \mid \varepsilon,
\]
\[
B\to b.
\]

Тогда:
\[
FIRST(A)=\{a,\varepsilon\},
\]
\[
FIRST(B)=\{b\},
\]
\[
FIRST(S)=\{a,b\}.
\]

Поскольку после $A$ стоит $B$:
\[
FOLLOW(A)=\{b\}.
\]

Поскольку $S$ --- стартовый символ:
\[
FOLLOW(S)=\{EOF\}.
\]

Поскольку $B$ стоит в конце продукции $S\to AB$:
\[
FOLLOW(B)=\{EOF\}.
\]

\subsection{LL(1)-анализ}

LL(1)-парсер:
\begin{itemize}
    \item читает вход слева направо;
    \item строит левосторонний вывод;
    \item использует один символ предпросмотра.
\end{itemize}

\[
LL(1)=
\text{Left-to-right scan, Leftmost derivation, 1 lookahead}.
\]

\subsubsection*{Условие LL(1)}

Для каждого нетерминала $A$ и любых двух различных продукций:
\[
A\to \alpha,
\]
\[
A\to \beta
\]
должны выполняться условия:
\[
FIRST(\alpha)\cap FIRST(\beta)=\varnothing,
\]
а если
\[
\varepsilon\in FIRST(\alpha),
\]
то
\[
FIRST(\beta)\cap FOLLOW(A)=\varnothing.
\]

Аналогично, если
\[
\varepsilon\in FIRST(\beta),
\]
то
\[
FIRST(\alpha)\cap FOLLOW(A)=\varnothing.
\]

\subsection{Построение таблицы LL(1)}

Для каждой продукции:
\[
A\to \alpha
\]
выполняются шаги:

\begin{enumerate}
    \item Для каждого терминала
    \[
    a\in FIRST(\alpha)\setminus\{\varepsilon\}
    \]
    занести продукцию $A\to\alpha$ в ячейку
    \[
    M[A,a].
    \]
    \item Если
    \[
    \varepsilon\in FIRST(\alpha),
    \]
    то для каждого
    \[
    b\in FOLLOW(A)
    \]
    занести продукцию $A\to\alpha$ в ячейку
    \[
    M[A,b].
    \]
    \item Если в какую-либо ячейку попадает более одной продукции, грамматика не является LL(1).
\end{enumerate}

\subsubsection*{Пример}

Грамматика:
\[
S\to aS \mid b.
\]

\[
FIRST(aS)=\{a\},
\]
\[
FIRST(b)=\{b\}.
\]

Таблица:
\[
M[S,a]=S\to aS,
\]
\[
M[S,b]=S\to b.
\]

Разбор строки:
\[
aaab.
\]

Стек и вход:

\[
\begin{array}{c|c|c}
\text{Стек} & \text{Вход} & \text{Действие}\\
\hline
SEOF & aaabEOF & S\to aS\\
aSEOF & aaabEOF & \text{съесть } a\\
SEOF & aabEOF & S\to aS\\
aSEOF & aabEOF & \text{съесть } a\\
SEOF & abEOF & S\to aS\\
aSEOF & abEOF & \text{съесть } a\\
SEOF & bEOF & S\to b\\
bEOF & bEOF & \text{съесть } b\\
EOF & EOF & \text{принять}
\end{array}
\]

\subsection{Левая рекурсия}

\begin{definition}
Грамматика имеет \textbf{левую рекурсию}, если существует нетерминал $A$ такой, что
\[
A\Rightarrow^+ A\alpha.
\]
\end{definition}

Непосредственная левая рекурсия:
\[
A\to A\alpha \mid \beta.
\]

LL-парсеры не могут напрямую работать с левой рекурсией, так как рекурсивный спуск зациклится.

Преобразование:
\[
A\to A\alpha_1 \mid \cdots \mid A\alpha_m \mid \beta_1 \mid \cdots \mid \beta_n
\]
заменяется на:
\[
A\to \beta_1A' \mid \cdots \mid \beta_nA',
\]
\[
A'\to \alpha_1A' \mid \cdots \mid \alpha_mA' \mid \varepsilon.
\]

\subsubsection*{Пример}

Было:
\[
E\to E+T \mid T.
\]

После устранения левой рекурсии:
\[
E\to TE',
\]
\[
E'\to +TE' \mid \varepsilon.
\]

\subsection{Левая факторизация}

Левая факторизация устраняет общий префикс альтернатив.

Если:
\[
A\to \alpha\beta_1 \mid \alpha\beta_2,
\]
то заменяем на:
\[
A\to \alpha A',
\]
\[
A'\to \beta_1 \mid \beta_2.
\]

Пример:
\[
S\to if\ E\ then\ S\ else\ S \mid if\ E\ then\ S.
\]

После факторизации:
\[
S\to if\ E\ then\ S\ S',
\]
\[
S'\to else\ S \mid \varepsilon.
\]

\subsection{LR-анализ}

LR-парсер:
\begin{itemize}
    \item читает вход слева направо;
    \item строит правосторонний вывод в обратном порядке;
    \item использует стек состояний;
    \item работает с действиями shift и reduce.
\end{itemize}

\[
LR(k)=
\text{Left-to-right scan, Rightmost derivation in reverse, }k\text{ lookahead}.
\]

\subsection{Расширенная грамматика}

Для построения LR-анализатора грамматику расширяют новой стартовой продукцией:
\[
S'\to S.
\]

Цель:
\[
S'\to S\cdot
\]
означает успешное завершение разбора.

\subsection{LR(0)-элементы}

\begin{definition}
\textbf{LR(0)-элемент} --- продукция грамматики с точкой в правой части.
\end{definition}

Если есть продукция:
\[
A\to XYZ,
\]
то возможные LR(0)-элементы:
\[
A\to \cdot XYZ,
\]
\[
A\to X\cdot YZ,
\]
\[
A\to XY\cdot Z,
\]
\[
A\to XYZ\cdot.
\]

Точка показывает, какая часть правой части уже распознана.

\subsection{Операция closure}

Если в множестве элементов есть:
\[
A\to \alpha\cdot B\beta,
\]
то для каждой продукции:
\[
B\to \gamma
\]
нужно добавить:
\[
B\to \cdot\gamma.
\]

Алгоритм:
\begin{enumerate}
    \item Начать с заданного множества элементов $I$.
    \item Если в $I$ есть элемент вида
    \[
    A\to \alpha\cdot B\beta,
    \]
    добавить все элементы
    \[
    B\to \cdot\gamma.
    \]
    \item Повторять, пока множество не перестанет изменяться.
\end{enumerate}

\subsection{Операция goto}

Для множества элементов $I$ и грамматического символа $X$:
\[
goto(I,X)=closure(\{A\to \alpha X\cdot \beta \mid A\to \alpha\cdot X\beta\in I\}).
\]

\subsection{Каноническая коллекция LR(0)-множеств}

Алгоритм:
\begin{enumerate}
    \item Построить начальное множество:
    \[
    I_0=closure(\{S'\to \cdot S\}).
    \]
    \item Для каждого построенного множества $I$ и каждого символа $X$ вычислить:
    \[
    goto(I,X).
    \]
    \item Если результат непустой и ещё не встречался, добавить его.
    \item Повторять до неподвижной точки.
\end{enumerate}

\subsection{Пример LR(0)-построения}

Грамматика:
\[
S'\to S,
\]
\[
S\to a.
\]

Начальное множество:
\[
I_0=closure(\{S'\to \cdot S\}).
\]

Поскольку после точки стоит $S$, добавляем:
\[
S\to \cdot a.
\]

Итак:
\[
I_0=
\{
S'\to \cdot S,\
S\to \cdot a
\}.
\]

Переход по $S$:
\[
goto(I_0,S)=\{S'\to S\cdot\}=I_1.
\]

Переход по $a$:
\[
goto(I_0,a)=\{S\to a\cdot\}=I_2.
\]

Действия:
\[
I_0 \xrightarrow{a} I_2,
\]
\[
I_0 \xrightarrow{S} I_1.
\]

В состоянии $I_2$ выполняется reduce по продукции:
\[
S\to a.
\]

В состоянии $I_1$ выполняется accept.

\subsection{Действия LR-парсера}

LR-парсер использует таблицы:
\[
ACTION[state, terminal],
\]
\[
GOTO[state, nonterminal].
\]

Возможные действия:

\begin{itemize}
    \item \textbf{shift} $s$ --- считать текущий токен и перейти в состояние $s$;
    \item \textbf{reduce} $A\to\alpha$ --- заменить на стеке правую часть $\alpha$ нетерминалом $A$;
    \item \textbf{accept} --- успешное завершение;
    \item \textbf{error} --- синтаксическая ошибка.
\end{itemize}

\subsubsection*{Алгоритм LR-разбора}

Вход:
\[
wEOF.
\]

Стек содержит состояния.

\begin{enumerate}
    \item Поместить в стек начальное состояние $0$.
    \item Пусть $s$ --- верхнее состояние стека, $a$ --- текущий входной токен.
    \item Посмотреть:
    \[
    ACTION[s,a].
    \]
    \item Если действие равно $shift\ t$, то:
    \begin{enumerate}
        \item поместить $a$ и состояние $t$ в стек;
        \item перейти к следующему входному токену.
    \end{enumerate}
    \item Если действие равно $reduce\ A\to\beta$, то:
    \begin{enumerate}
        \item удалить со стека $|\beta|$ грамматических символов и состояний;
        \item пусть $s'$ --- новое верхнее состояние;
        \item поместить $A$;
        \item поместить состояние
        \[
        GOTO[s',A].
        \]
    \end{enumerate}
    \item Если действие равно $accept$, принять вход.
    \item Если действие равно $error$, сообщить об ошибке.
\end{enumerate}

\subsection{Конфликты LR-анализа}

При построении LR-таблиц могут возникнуть конфликты.

\subsubsection*{Shift/reduce-конфликт}

Возникает, когда в одной ячейке таблицы возможны:
\[
shift
\]
и
\[
reduce.
\]

Классический пример --- висячий \texttt{else}:
\[
S\to if\ E\ then\ S,
\]
\[
S\to if\ E\ then\ S\ else\ S.
\]

Обычно конфликт разрешается в пользу \texttt{shift}, чтобы \texttt{else} связывался с ближайшим \texttt{if}.

\subsubsection*{Reduce/reduce-конфликт}

Возникает, когда в одной ячейке возможны два разных reduce.

Например:
\[
A\to x,
\]
\[
B\to x.
\]

После чтения $x$ парсер может не знать, сворачивать ли $x$ в $A$ или в $B$.

\subsection{SLR, LR(1), LALR}

\subsubsection*{SLR(1)}

SLR использует LR(0)-элементы и множества FOLLOW.

Если в состоянии есть элемент:
\[
A\to \alpha\cdot,
\]
то reduce по продукции $A\to\alpha$ помещается только для терминалов:
\[
a\in FOLLOW(A).
\]

\subsubsection*{LR(1)}

LR(1)-элемент имеет вид:
\[
[A\to \alpha\cdot\beta,\ a],
\]
где $a$ --- символ предпросмотра.

Смысл: продукцию $A\to\alpha\beta$ можно использовать в контексте, где после $A$ ожидается $a$.

LR(1) точнее SLR, но обычно порождает больше состояний.

\subsubsection*{LALR(1)}

LALR(1) строится из LR(1), но состояния с одинаковым LR(0)-ядром объединяются.

\[
\text{LALR(1)}:
\]
\begin{itemize}
    \item мощнее SLR(1);
    \item обычно компактнее LR(1);
    \item используется в классических генераторах вроде yacc и bison.
\end{itemize}

\subsection{Автоматическая генерация анализаторов}

\subsubsection{Lex}

\textbf{lex} --- генератор лексических анализаторов.

Файл lex обычно имеет структуру:
\[
\begin{array}{c}
\%\{\\
\text{C-код объявлений}\\
\%\}\\
\%\%\\
\text{лексические правила}\\
\%\%\\
\text{дополнительный C-код}
\end{array}
\]

Минимальный пример:

\begin{verbatim}
%{
#include "y.tab.h"
%}

%%
[0-9]+                 { return NUM; }
[a-zA-Z_][a-zA-Z0-9_]* { return ID; }
"+"                    { return PLUS; }
[ \t\n]+               { /* skip whitespace */ }
.                      { return yytext[0]; }
%%

int yywrap(void) {
    return 1;
}
\end{verbatim}

Правило:
\[
\text{\texttt{[0-9]+}}
\]
распознаёт целые числа.

Переменная \texttt{yytext} содержит текст распознанной лексемы.

Функция \texttt{yylex()} возвращает следующий токен.

\subsubsection{Yacc}

\textbf{yacc} --- генератор синтаксических анализаторов, обычно использующий LALR(1)-анализ.

Файл yacc имеет структуру:
\[
\begin{array}{c}
\%\{\\
\text{C-код объявлений}\\
\%\}\\
\text{объявления токенов}\\
\%\%\\
\text{грамматические правила}\\
\%\%\\
\text{дополнительный C-код}
\end{array}
\]

Минимальный пример калькулятора:

\begin{verbatim}
%token NUM
%left '+'
%left '*'

%%
expr:
      expr '+' expr
    | expr '*' expr
    | '(' expr ')'
    | NUM
    ;
%%
\end{verbatim}

Директивы:
\[
\%left '+'
\]
и
\[
\%left '*'
\]
задают левую ассоциативность.

Порядок объявлений приоритетов обычно идёт от меньшего приоритета к большему. Поэтому \texttt{*}, объявленный после \texttt{+}, имеет больший приоритет.

\subsubsection*{Связка lex и yacc}

Обычно:
\begin{itemize}
    \item lex генерирует функцию \texttt{yylex()};
    \item yacc генерирует функцию \texttt{yyparse()};
    \item \texttt{yyparse()} вызывает \texttt{yylex()} для получения токенов.
\end{itemize}

Схема:
\[
\texttt{input}
\to
\texttt{yylex()}
\to
\texttt{tokens}
\to
\texttt{yyparse()}.
\]

\subsection{ANTLR}

\textbf{ANTLR} --- современный генератор лексических и синтаксических анализаторов.

В отличие от классической пары lex/yacc, ANTLR обычно позволяет описывать лексические и синтаксические правила в одном файле грамматики.

Пример грамматики:

\begin{verbatim}
grammar Expr;

prog: expr EOF;

expr
    : expr '*' expr
    | expr '+' expr
    | INT
    | '(' expr ')'
    ;

INT: [0-9]+;
WS : [ \t\r\n]+ -> skip;
\end{verbatim}

ANTLR поддерживает генерацию парсеров для разных целевых языков.

Типичная структура:
\[
\text{ANTLR-грамматика}
\to
\text{лексер}
+
\text{парсер}
+
\text{visitor/listener}.
\]

\subsubsection*{Listener и Visitor}

ANTLR может генерировать два основных интерфейса обхода дерева:

\begin{itemize}
    \item \textbf{Listener} --- автоматически вызывается при входе и выходе из узлов дерева;
    \item \textbf{Visitor} --- пользователь явно управляет обходом дерева.
\end{itemize}

Пример visitor-идеи:
\[
visitExpr(node)
\]
возвращает значение выражения или строит фрагмент AST.

\subsection{Парсер-комбинаторы}

\textbf{Парсер-комбинаторы} --- подход, при котором парсер является обычной функцией, а сложные парсеры строятся из простых с помощью комбинаторов.

Типовая абстракция:
\[
Parser\ a = String \to [(a,String)].
\]

То есть парсер значения типа $a$ принимает строку и возвращает список возможных результатов:
\[
\text{результат разбора},\ \text{остаток входа}.
\]

\subsubsection*{Базовые парсеры}

Парсер одного символа:

\[
item : Parser\ Char.
\]

Он читает первый символ входа.

Парсер, проверяющий предикат:

\[
satisfy(p)
\]
читает символ, если он удовлетворяет предикату $p$.

Парсер конкретного символа:
\[
char(c)=satisfy(\lambda x.\ x=c).
\]

\subsubsection*{Основные комбинаторы}

\[
p \mid q
\]
--- альтернатива: попробовать $p$, если не получилось --- $q$.

\[
p \gg= f
\]
--- последовательная композиция: выполнить $p$, затем по его результату выбрать следующий парсер.

\[
many(p)
\]
--- повторить $p$ ноль или более раз.

\[
many1(p)
\]
--- повторить $p$ один или более раз.

\subsubsection*{Пример}

Парсер натурального числа:

\begin{verbatim}
digit  = satisfy isDigit
number = many1 digit
\end{verbatim}

Парсер суммы чисел:

\begin{verbatim}
expr = number `chainl1` plus
plus = char '+' >> return (+)
\end{verbatim}

Идея \texttt{chainl1}: разобрать цепочку левоассоциативных операций.

Строка:
\[
\texttt{1+2+3}
\]
разбирается как:
\[
(1+2)+3.
\]

\subsection{Сравнение подходов}

\[
\begin{array}{|c|c|c|c|}
\hline
\text{Инструмент} & \text{Описание лексики} & \text{Описание синтаксиса} & \text{Типичный метод}\\
\hline
lex & \text{регулярные выражения} & - & \text{ДКА}\\
\hline
yacc & - & \text{КС-грамматика} & \text{LALR(1)}\\
\hline
ANTLR & \text{лексические правила} & \text{грамматические правила} & \text{адаптивный LL(*)}\\
\hline
\text{парсер-комбинаторы} & \text{код} & \text{код} & \text{рекурсивный спуск}\\
\hline
\end{array}
\]

\subsection{Преимущества и недостатки}

\subsubsection*{lex/yacc}

Преимущества:
\begin{itemize}
    \item классический и хорошо изученный подход;
    \item эффективные анализаторы;
    \item хорошая теоретическая база;
    \item удобно для языков с LALR-грамматиками.
\end{itemize}

Недостатки:
\begin{itemize}
    \item не всегда удобная диагностика ошибок;
    \item конфликты shift/reduce и reduce/reduce могут быть трудны для начинающих;
    \item синтаксис инструментов исторически ориентирован на C.
\end{itemize}

\subsubsection*{ANTLR}

Преимущества:
\begin{itemize}
    \item единый файл грамматики для лексики и синтаксиса;
    \item генерация parse tree;
    \item listener/visitor-модель;
    \item поддержка нескольких целевых языков.
\end{itemize}

Недостатки:
\begin{itemize}
    \item грамматики требуют знания особенностей ANTLR;
    \item некоторые неоднозначности и приоритеты нужно описывать явно;
    \item для очень низкоуровневой оптимизации может быть менее прозрачен, чем рукописный парсер.
\end{itemize}

\subsubsection*{Парсер-комбинаторы}

Преимущества:
\begin{itemize}
    \item парсер пишется на основном языке программирования;
    \item удобно строить AST;
    \item высокая композиционность;
    \item хорошие возможности для встроенных DSL.
\end{itemize}

Недостатки:
\begin{itemize}
    \item наивные реализации могут быть медленными;
    \item левая рекурсия обычно приводит к зацикливанию;
    \item диагностика ошибок требует отдельной работы;
    \item производительность зависит от реализации.
\end{itemize}

\subsection{Минимальный сквозной пример}

Рассмотрим язык арифметических выражений:
\[
1+2*3.
\]

Лексические правила:
\[
\text{\texttt{[0-9]+}} \to NUM,
\]
\[
+ \to PLUS,
\]
\[
* \to MUL.
\]

Поток токенов:
\[
\langle NUM(1), PLUS, NUM(2), MUL, NUM(3), EOF \rangle.
\]

Грамматика:
\[
E\to E+T \mid T,
\]
\[
T\to T*F \mid F,
\]
\[
F\to NUM \mid (E).
\]

Разбор:
\[
\begin{aligned}
E &\Rightarrow E+T \\
  &\Rightarrow T+T \\
  &\Rightarrow F+T \\
  &\Rightarrow NUM+T \\
  &\Rightarrow NUM+T*F \\
  &\Rightarrow NUM+F*F \\
  &\Rightarrow NUM+NUM*NUM.
\end{aligned}
\]

AST:
\[
+
\left(
1,
*
\left(
2,
3
\right)
\right).
\]

То есть выражение интерпретируется как:
\[
1+(2*3),
\]
а не как:
\[
(1+2)*3.
\]

\subsection{Итог}

Автоматическое построение лексических и синтаксических анализаторов основано на строгом соответствии между формальными описаниями и алгоритмическими моделями:

\[
\text{регулярные выражения}
\leftrightarrow
\text{конечные автоматы}
\leftrightarrow
\text{лексеры},
\]

\[
\text{контекстно-свободные грамматики}
\leftrightarrow
\text{LL/LR-анализ}
\leftrightarrow
\text{парсеры}.
\]

Инструменты вроде lex и yacc реализуют классическую схему:
\[
\text{регулярные выражения} + \text{КС-грамматика}
\to
\text{лексер} + \text{LALR-парсер}.
\]

ANTLR предоставляет более современную интегрированную модель генерации анализаторов, а парсер-комбинаторы позволяют описывать парсеры как обычные программы, сохраняя математическую композиционность.

\end{document}
